\documentclass[11pt]{article}
\usepackage{series}
\usepackage{amsmath,amssymb,amsthm}
\usepackage{geometry}
\usepackage{hyperref}

\geometry{margin=1in}

\title{\textbf{Particles and Forces as Coherence Basins\\ and Capacity Gradients}}
\author{Peter Nero}
\date{January 2026}

\begin{document}
\maketitle

\begin{abstract}
We show how particles and forces arise as effective structures in admissible descriptions
built from coherence basins and coherence-capacity transport, under the same stability,
regularity, and closure assumptions used elsewhere in the coherence-capacity framework. Building on the transport theory developed in
``Dynamics of Coherence Capacity: Transport, Concentration, and Exhaustion,''
we show that localized, persistent excitations correspond to stable
coherence basins, while forces arise from gradients and flows of coherence
capacity. Newtonian motion, the Lorentz force, and confinement phenomena
are obtained as basin-centroid and coherence-transport effects without
introducing fundamental force carriers. Particle number nonconservation,
charge conservation, and universality of gravity follow naturally from
the basin framework.
\end{abstract}

\section{Introduction}

In the preceding papers we identified coherence capacity as the fundamental
resource governing effective physical descriptions and developed its
transport, concentration, and exhaustion dynamics. In this work we show
how the familiar ontology of particles and forces arises from that same
structure.

The central claim of this paper is simple:
\emph{particles are stable coherence basins, and forces are the shadows of
coherence-capacity gradients acting on those basins}.

This viewpoint eliminates the need to treat particles as fundamental
point objects or forces as primitive interactions. Instead, particles and
forces emerge from the same projection-based stability constraints that
govern gravity, time, and irreversibility.

\section{Effective States and Coherence Basins}

\subsection{Effective state space}

Let $Y$ denote the space of effective (observable) states obtained by
projection from the fundamental configuration space. Within admissible
regions, effective dynamics is well-defined and predictive.
\paragraph{Effective description level.}
The effective state space $Y$ should be understood as the image of an admissible projection
from the coherent sector of the underlying state space, equipped with whatever minimal
structure (topological, measure-theoretic, or differentiable) is required to support the
effective equations used below. No claim is made that $Y$ exists globally or independently of
the admissible regime.

\subsection{Coherence basins}

\begin{definition}[Coherence basin]
A coherence basin $B_\alpha \subset Y$ is a connected region of effective
state space such that:
\begin{itemize}
\item projection stability holds throughout $B_\alpha$,
\item trajectories entering $B_\alpha$ remain in $B_\alpha$ under
admissible evolution,
\item disturbances within $B_\alpha$ are contractively damped.
\end{itemize}
\end{definition}

Coherence basins are the effective analogues of attractors or fixed-point
regions in dynamical systems.

\begin{definition}[Particle]
A particle is a localized, persistent coherence basin whose support in
effective space remains bounded and identifiable over admissible
evolution.
\end{definition}

This definition makes no reference to point-like ontology or fundamental
fields. Persistence and localization are properties of basin stability.

\section{Basin Density and Centroid Dynamics}
\subsection{Effective density representation (assumption)}

In admissible regimes where a smooth effective chart exists, we assume that the restriction
of the coherence basin admits an absolutely continuous representation with respect to the
effective spatial measure, yielding a basin density $\rho_\alpha(x,t)$. This representation
is effective and slab-local; it does not assert a fundamental density or fluid ontology.
All subsequent continuum equations are understood in this restricted, encoding-level sense.
\begin{lemma}[Preservation of basin identity under effective representation]
  Within admissible regimes, the effective density representation $\rho_\alpha(x,t)$ uniquely
tracks the identity of the underlying coherence basin. Distinct coherence basins cannot yield
identical effective density evolutions on overlapping admissible charts, except at
admissibility barriers where basin merge--split occurs.  
\end{lemma}


\subsection{Basin density}

To describe motion of a coherence basin, we introduce a basin density
$\rho_\alpha(x,t)$ on an effective spatial slice $\Sigma$.

Normalization:
\[
M_\alpha := \int_\Sigma \rho_\alpha(x,t)\,d^3x
\]
is conserved within admissible regimes.

\subsection{Centroid definition}

\begin{definition}[Basin centroid]
The centroid of a coherence basin is
\[
X^i(t) := \frac{1}{M_\alpha}\int_\Sigma x^i\,\rho_\alpha(x,t)\,d^3x.
\]
\end{definition}

The centroid is the effective worldline of the particle.

\section{Continuity and Momentum Balance}

\subsection{Continuity equation}

Within an admissible basin, the reduced dynamics yields a continuity
equation:
\[
\partial_t \rho_\alpha + \nabla\cdot(\rho_\alpha v) = 0,
\]
where $v$ is the effective velocity field induced by the underlying
dynamics and projection.

\paragraph{Constitutive closure.}
The continuity and momentum-balance equations below are constitutive relations of the
effective description. Their existence follows from the assumed regularity of the basin
density representation and does not introduce new microscopic dynamics.


\subsection{Momentum balance}

Assume the existence of an effective momentum flux tensor $\Pi^{ij}$ such
that:
\[
\partial_t(\rho_\alpha v^i) + \partial_j \Pi^{ij}
= \rho_\alpha a^i_{\mathrm{eff}},
\]
where $a^i_{\mathrm{eff}}$ encodes bias induced by coherence-capacity
gradients and internal connections.

\section{Newtonian Motion from Basin Centroids}

\begin{theorem}[Centroid equation of motion, effective]
If the basin is narrow compared to the scale of variation of
$a^i_{\mathrm{eff}}$, then the centroid satisfies
\[
M_\alpha \ddot X^i = F^i(X),
\]
where
\[
F^i(X) := M_\alpha a^i_{\mathrm{eff}}(X).
\]
\end{theorem}

\begin{proof}
The result holds in the narrow-basin, slowly varying limit of the effective continuum
representation.
Integrating the momentum balance equation over space and using localization
of $\rho_\alpha$ yields the stated result.
\end{proof}

\begin{corollary}
Newton's second law arises as the narrow-basin limit of coherence-basin
dynamics.
\end{corollary}

This law is not fundamental but emergent: the ``force'' reflects bias in
capacity transport rather than interaction between point objects.

\section{Interpretation}

The derivation above shows that classical particle motion emerges whenever
coherence basins are sufficiently localized and capacity gradients vary
slowly. Mass corresponds to the integrated coherence weight $M_\alpha$ of
the basin; force corresponds to spatial variation of the effective
capacity-induced acceleration field. This quantity encodes inertial resistance of the basin to chart reconfiguration, not a
fundamental mass density; its appearance as an integral reflects the chosen effective
representation.


No assumption of fundamental forces has been made.

\section{Electromagnetism as Phase-Coherence Transport}

\subsection{Phase coherence and internal connections}

In addition to spatial localization, coherence basins generally carry
internal phase structure inherited from the projection of the fundamental
configuration space. We model this by associating to each basin a complex
amplitude
\[
\psi(x) = \sqrt{\rho(x)}\,e^{i\theta(x)},
\]
where $\rho$ is the basin density and $\theta$ is an internal phase.

Gauge freedom arises because only relative phase is physically meaningful.
Accordingly, the effective description involves a connection $A_\mu(x)$,
and physical phase gradients appear only in the gauge-covariant
combination
\[
p_\mu := \partial_\mu \theta - q A_\mu,
\]
where $q$ labels the basin's coupling to the internal phase connection.

\subsection{Phase-gradient transport}

Assuming the Hamilton--Jacobi constraint obtained in the eikonal limit of coherent wave
reconstruction, within an admissible basin, coherent evolution preserves the Hamilton--Jacobi
constraint
\[
g^{\mu\nu} p_\mu p_\nu = m^2,
\]
which may be viewed as the eikonal limit of coherent wave dynamics.

Taking a covariant derivative along the basin centroid worldline
$X^\mu(\tau)$ yields
\[
\frac{D p_\mu}{D\tau}
= -q\,\partial_\mu A_\nu\,\dot X^\nu
+ q\,\partial_\nu A_\mu\,\dot X^\nu.
\]

Introducing the field strength
\[
F_{\mu\nu} := \partial_\mu A_\nu - \partial_\nu A_\mu,
\]
we obtain the effective equation of motion
\[
m\,\frac{D \dot X^\mu}{D\tau} = q\,F^\mu{}_{\nu}\,\dot X^\nu.
\]

\begin{theorem}[Lorentz force as effective coherence-transport law]
The Lorentz force arises as the condition that phase coherence be preserved
under admissible transport of coherence basins.
\end{theorem}
This result holds at the effective description level and presupposes the phase-coherent
transport structure established in the underlying reconstruction.

Electromagnetism therefore does not represent a fundamental force carrier,
but the bookkeeping of phase-coherence preservation during basin motion.

\section{Gauge Charge and Conservation Laws}

\subsection{Charge as a basin invariant}

The quantity $q$ appearing in the covariant phase gradient is an invariant
label of the coherence basin. It characterizes how the basin transforms
under internal phase rotations.

\begin{definition}[Gauge charge]
Gauge charge is the representation label associated with internal
phase-coherence transport of a coherence basin.
\end{definition}

\subsection{Charge conservation}

Because admissible evolution preserves basin identity except at barriers,
gauge charge is conserved along admissible basin trajectories, except at admissibility
barriers.


\begin{theorem}[Charge conservation]
Within admissible regimes, gauge charge is conserved along coherence-basin
trajectories.
\end{theorem}

\begin{proof}
Charge labels are topological invariants of the basin structure and cannot
change under continuous admissible evolution.
\end{proof}

In contrast, particle number is not protected: basins may merge, split, or
dissolve at admissibility barriers, leading to particle creation and
annihilation.

\section{Confinement as Coherence Protection}

\subsection{Non-Abelian coherence strain}

In sectors with non-Abelian internal structure, coherence transport imposes
additional strain. Maintaining non-singlet coherence across extended
regions requires preserving correlated internal orientations, which
rapidly consumes coherence capacity.

\subsection{Area-law cost}

Consider separating two coherence basins carrying non-Abelian charge.
Maintaining admissibility requires preserving a coherent ``string'' of
internal alignment between them. The capacity cost of this configuration
grows with the area swept by the separation process.

\begin{theorem}[Confinement from capacity exhaustion, conditional]
Assume bounded transport flux and persistent coherence strain under non-Abelian separation.
Non-singlet coherence basins in non-Abelian sectors experience capacity
exhaustion under separation, forcing recombination into singlet basins.
\end{theorem}

\begin{proof}
Persistent capacity strain along the separation direction induces focusing
of capacity flux and eventual formation of a bottleneck. The only
admissible configuration that avoids capacity exhaustion is one in which
net non-Abelian charge is neutralized locally.
\end{proof}

This reproduces confinement and the area-law behavior of Wilson loops as
consequences of coherence protection rather than fundamental interaction
potentials.

\section{Summary of Particle and Force Emergence}

We summarize the emergent picture:

\begin{itemize}
\item Particles are stable coherence basins.
\item Mass is the integrated coherence weight of a basin.
\item Motion arises from capacity-gradient bias.
\item Gravity appears as the universal response to total coherence-capacity distribution in
effective geometric encodings.
\item Electromagnetism preserves phase coherence under transport.
\item Gauge charges are basin invariants.
\item Confinement enforces coherence protection in non-Abelian sectors.
\end{itemize}

No fundamental point particles or force carriers are required.

\section{Discussion}

Particles and forces emerge here as structural features of projection-based
effective descriptions with finite coherence capacity. This framework
explains why gravity is universal, why gauge interactions are quantized,
and why particle number is not conserved, while charge is.

The derivations rely only on stability, projection, and capacity transport,
and therefore apply broadly to any effective theory with similar
structure.

\section{Conclusion}

We have shown that particles and forces arise naturally from the dynamics
of coherence capacity. Stable coherence basins behave as particles, while
gradients and transport of coherence capacity generate the effective forces
governing their motion.

This completes the coherence-capacity trilogy at the level of effective emergence:
\begin{itemize}
\item coherence capacity as the fundamental resource,
\item its transport, concentration, and exhaustion,
\item and the emergence of particles and forces.
\end{itemize}

Further work will address cosmology, entropy bounds, and quantum field
theory reconstruction within this framework.

\end{document}

\end{document}
